\documentclass[preprint,12pt]{elsarticle}
% Supplementary material of manuscript/rebuild/main.tex (Brief 29). Its tables are written by
% src/make_rebuild_registry.py and src/make_rebuild_tables.py from the results files; its numbers are traced by
% manuscript/numbers_check.py --rebuild --supplement.
\usepackage[T1]{fontenc}
\usepackage{lmodern}
\usepackage{amsmath}
\usepackage{graphicx}
\usepackage{booktabs}
\usepackage{longtable}
\usepackage{array}
\usepackage{textcomp}
\usepackage{xurl}
\usepackage[hidelinks]{hyperref}
\newcommand{\degC}{\,\textdegree C}
\renewcommand{\thesection}{S\arabic{section}}
\renewcommand{\thetable}{S\arabic{table}}
\renewcommand{\thefigure}{S\arabic{figure}}
\biboptions{sort&compress}
\setlength{\emergencystretch}{1.5em}
\makeatletter
\def\ps@pprintTitle{\let\@oddhead\@empty\let\@evenhead\@empty\def\@oddfoot{\reset@font\hfil\thepage\hfil}\let\@evenfoot\@oddfoot}
\makeatother

\begin{document}

\begin{frontmatter}
\title{Supplementary material for: The heat a neighbouring lithium-ion cell needs for thermal runaway propagation depends on how the cell is modelled}
\author{Siddharth Satte}
\address{Independent researcher, Pune, Maharashtra, India}
\end{frontmatter}

This supplement holds the record of the study (its fifteen steps and every prediction registered before their runs), the numbers behind the figures of the main text test by test, the checks, the inputs of the bounds on what each path can deliver, and how the one published neighbour history was read. Everything in it can be read without the archive; the archive~\cite{satte2026repo} holds the code, the run records and a log of every decision, for anyone who wants to repeat a result.

\section{The labels used in the registry}
\label{sec:s-key}

The study's records, and so the predictions in Table~\ref{tab:s-registry}, use short labels. In plain words:
\begin{itemize}
\item Trigger, neighbour: the failing cell, modelled as one temperature, and the cell beside it. A lumped neighbour is one temperature; a resolved neighbour is a three-dimensional temperature field.
\item Hottest node, hot spot: the resolved neighbour's hottest point at each time. Mean: the resolved neighbour's heat-capacity-weighted mean temperature. Far face: the resolved neighbour's surface point at mid-length opposite the line of contact.
\item C1, the local criterion: the hot spot reaches 255.7\degC. C2, the whole-cell clock: the lumped neighbour reaches 255.7\degC. C3: the lumped neighbour reaches 150\degC. C4 and C5: the separator's melting temperature, 138.5\degC, at the hot spot or in the lumped neighbour. C6 and C7: C1 and C4 with the SEI already consumed. C8: the resolved neighbour's mean reaches 255.7\degC{} within 900~s, the cell's own runaway.
\item $E_\mathrm{b}$: the trigger's body energy; $E_\mathrm{e}=3.105\,E_\mathrm{b}$: the energy its ejecta carry, with 3.105 the median ratio of ejected to body energy in the ten nail tests of this cell in the Battery Failure Databank~\cite{finegan2024databank}; $E_\mathrm{fix}=76.74$~kJ: the Databank's median ejected energy, one for every test.
\item $f$: the fraction of $E_\mathrm{e}$ delivered to the neighbour through the half of its can facing the trigger. $f^*$ and $E^*$: a criterion's threshold, the least fraction or energy at which it is met. $f_\mathrm{t}$: the fraction at which a criterion is met at a propagated test's measured time. $f_\mathrm{max}$: the fraction at which the far face peaks at a non-propagated test's measured maximum.
\item Low, central and high contact: a contact conductance of 200, 2,000 and 3,700~W\,m$^{-2}$\,K$^{-1}$ between the cans.
\item A1 and A2 in the predictions of step 13: the skin model and the cross-section conduction model; B1 and B2: the ignition temperature and the two-body model (Section~2.5 and Appendix~B of the main text). In step 14, A1 to A6 are the paths of Table~\ref{tab:s-inputs}, in its order.
\end{itemize}

\section{The fifteen steps}
\label{sec:s-steps}

The study ran in fifteen steps (Table~\ref{tab:s-steps}). Each step's predictions file was committed and pushed to the public repository before any run of the step's registered list, with the outcome expected for each prediction, and was not changed afterwards; the repository records each file's commit, the time of its push and the hash of each step's run list~\cite{satte2026repo}. Two departures are recorded there. While the eleventh step's code was tested, one full run under a kinetic draw was made before its predictions, and its result is stated in them; and while that step's runs went, the scoring was once run on 178 of its 216 records to test a figure's script, and the verdicts it printed were seen, after which no code, rule or run changed.

\IfFileExists{tables/s1_steps.tex}{\begin{table}[!htbp]
\centering
\scriptsize
\caption{The study's fifteen steps, each with its own predictions committed to the public repository and pushed before its runs (the date is that of the push, UTC), and how many held and failed; a failure is foreseen when the predictions file recorded it as expected before the runs.}
\label{tab:s-steps}
\setlength{\tabcolsep}{3pt}
\begin{tabular}{rp{0.44\textwidth}lrrrr}
\toprule
Step & What it asked & Registered & Predictions & Held & Failed & Foreseen \\
\midrule
1 & Passive kinetics along the archived histories of the earlier two-cell study: when the chemistry starts at the hot spot, the mean and the lumped node & 2026-10-10 & 4 & 4 & 0 & 0 \\
2 & The reaction heat fed back, one way, to both neighbours & 2026-10-10 & 3 & 2 & 1 & 0 \\
3 & The onsets over 200 kinetic draws and 30 thermal draws & 2026-10-12 & 6 & 3 & 3 & 2 \\
4 & Which submodel of the kinetic posterior sets the lead of the hot spot & 2026-10-16 & 5 & 4 & 1 & 0 \\
5 & The cans in contact and joined by tabs, against the sixteen measured tests & 2026-10-21 & 7 & 1 & 6 & 0 \\
6 & Why the model does not propagate: the tests' 150\degC{} criterion and the trigger's measured durations & 2026-10-23 & 7 & 1 & 6 & 6 \\
7 & Seven ignition criteria scored test by test, each test with its own trigger & 2026-10-26 & 8 & 3 & 5 & 5 \\
8 & The delivered fraction of the ejected energy that each test's outcome requires & 2026-10-26 & 8 & 5 & 3 & 3 \\
9 & The delivered fraction read from each test's time and peak & 2026-10-29 & 8 & 5 & 3 & 3 \\
10 & The one published history, the thermocouple's angle and held-out scores & 2026-10-30 & 8 & 4 & 4 & 2 \\
11 & Bands of the central numbers under kinetic and thermal draws chosen by rank & 2026-11-02 & 6 & 2 & 4 & 1 \\
12 & The trigger held to calorimetry's body energy and one ejected energy for every test & 2026-11-03 & 4 & 4 & 0 & 0 \\
13 & The mechanism of each threshold: reduced models of skin conduction and of an ignition temperature & 2026-11-16 & 10 & 7 & 3 & 0 \\
14 & What each path can deliver to the neighbour, bounded from published inputs & 2026-11-17 & 6 & 4 & 2 & 2 \\
15 & The resolved neighbour's own runaway: a reactive cross-section and three-dimensional searches & 2026-11-18 & 8 & 7 & 1 & 0 \\
\midrule
All & & & 98 & 56 & 42 & 24 \\
\bottomrule
\end{tabular}
\end{table}
}{}

The main text reports the results of steps 5 to 15; steps 1 to 4 concern when the chemistry starts in an earlier geometry of the same cells, with a gap between them and busbars, and are summarised here only through their predictions. A failure counts as foreseen when the predictions file recorded it, before the runs, as expected, as fixed by earlier records or an energy balance, or as following from a condition it named that then held; a hedged ``may fail'' does not count. By that rule 24 of the 42 failures were foreseen.

\section{The registry of predictions}
\label{sec:s-registry}

\IfFileExists{tables/s2_registry.tex}{{\sloppy\scriptsize\setlength{\tabcolsep}{3pt}
\begin{longtable}{>{\raggedright\arraybackslash}p{0.05\textwidth}>{\raggedright\arraybackslash}p{0.065\textwidth}>{\raggedright\arraybackslash}p{0.425\textwidth}>{\raggedright\arraybackslash}p{0.24\textwidth}>{\raggedright\arraybackslash}p{0.1\textwidth}}
\caption{The 52 predictions whose runs the main text or this supplement reports, of the 98 the study registered: 34 held and 18 failed, 6 of the failures foreseen. Each is given as its predictions file states it (shortened where long), with the expectation recorded before the runs and the verdict. Labels are defined in Section~\ref{sec:s-key}.}\label{tab:s-registry}\\
\toprule
Step & & Prediction & Expectation & Verdict \\
\midrule
\endfirsthead
\multicolumn{5}{l}{\textit{Continued}} \\
\toprule
Step & & Prediction & Expectation & Verdict \\
\midrule
\endhead
5 & P1 & At the central contact conductance with the posterior-mean kinetics, the resolved neighbour propagates in series and parallel, its time to propagation within the measured 147 to 342 s. & Study's expectation: May fail: by the estimates above the hot spot stays below 255.7\degC, or, if the line of contact rises further, reaches it within the first minute. & failed \\
5 & P2 & In both connections, series and parallel, the lumped neighbour propagates more than 20 s later than the resolved one, or not at all within 900 s. & Study's expectation: Held, most likely as "not at all": its passive peak should lie near 145 to 170\degC. & held \\
5 & P3 & At the resolved model's propagation time, its mid-surface temperature lies between 150 and 255.7\,\textdegree C: the hot spot leads the surface a thermocouple sees. & Study's expectation: Toward the trigger, held if propagation occurs (that face is the line of contact); away from it, may fail below 150\degC. & failed \\
5 & P4 & The resolved model's mean mid-surface heating rate over the 60 s before propagation lies between 0.2 and 0.8\,\textdegree C/s (measured about 0.4). & Study's expectation: Uncertain. & failed \\
5 & P5 & The predicted series and parallel times to propagation differ by less than 30 s (the measured means differ by 8 s). & Study's expectation: Held if both propagate: the tabs carry 4 to 8 \% of the contact path. & failed \\
5 & P6 & The share of the 4,800 factorial runs in which the resolved neighbour propagates lies between 30 and 90 \% (measured 9 of 16, 56 \%). & Study's expectation: May fail low: the model lacks the vent, ejecta and flame heat, and the estimates above fall short of 255.7\degC{} at the central contact. & failed \\
5 & P7 & The contact and tab conductance range carries more of the spread in predicted time to propagation than the kinetic posterior (standard deviation ratio above 1), reversing step 3's onset finding. & Study's expectation: Held if enough runs propagate to give both spreads: the contact range changes the path 4.6-fold. & failed \\
7 & P1 & At the central contact with per-test triggers, C1 (hot spot at 255.7\degC) predicts propagation in fewer than 9 of the 16 tests. & Plan's expectation: Likely. Study's expectation: Holds, with eight. & held \\
8 & P1 & Of the nine propagated tests, $f^*$(C1) is 0 in the six step 7 hits and below 0.3 in the other three. & Plan's expectation: Likely. Study's expectation: Holds. & held \\
8 & P2 & Of the seven non-propagated tests, at least three, including the two step 7 false alarms at 0, have $f^*$(C1) below 0.3, so the propagated and non-propagated tests' $f^*$(C1) ranges overlap. & Plan's expectation: Likely. Study's expectation: Holds, with all seven below 0.3. & held \\
8 & P3 & $f^*$(C2) lies between 0.05 and 0.5 in at least twelve of the sixteen tests. & Plan's expectation: Uncertain. Study's expectation: Fails. & failed, foreseen \\
8 & P6 & Within each connection, $f^*$(C1) and $f^*$(C2) never increase with trigger energy (no inversion). & Plan's expectation: Likely; a check on the mapping. Study's expectation: Holds, with ties. & held \\
9 & P1 & With the 21.39 s delivery, C1 reproduces no propagated test's measured time within 30 s at any fraction at or above $f^*$'s lower end. & Plan's expectation: Likely. Study's expectation: Holds. & held \\
9 & P2 & With the 21.39 s delivery, C2 reproduces at least six of nine propagation times within 30 s, at a fraction between $f^*$'s lower end and 0.1 above its upper end. & Plan's expectation: Uncertain. Study's expectation: Holds. & held \\
9 & P3 & Every C2 time-matching fraction found in P2 lies no more than 0.02 above the upper end of $f^*$(C2). & Plan's expectation: Uncertain. Study's expectation: Uncertain; it turns on PP-Exp1. & held \\
9 & P4 & For the five low-maximum non-propagated tests, the peak-matching fraction at the thermocouple node lies below both $f^*$(C1) and $f^*$(C2). & Plan's expectation: Likely. Study's expectation: Fails, on ST-Exp1 and ST-Exp2, whatever the runs show. & failed, foreseen \\
9 & P5 & The peak-matching fractions (five non-propagated tests) and the C2 time-matching fractions (nine propagated tests) separate: the largest of the former is below the smallest of the latter. & Plan's expectation: Uncertain; this is the step's question. Study's expectation: Fails. & failed, foreseen \\
9 & P7 & The bracket on the ejected-to-body energy ratio rescales every time- and peak-matching fraction by 0.75 to 1.40 and changes P5's verdict in neither direction. & Plan's expectation: Likely. Study's expectation: Holds. & held \\
10 & P1 & Before propagation, the model under the 21.39 s delivery, at one end of its threshold bracket, fits PT-Exp4's measured history at the step 9 (far-face) node within 25 K root-mean-square. & Plan's expectation: uncertain. & held \\
10 & P6 & Held out, no predictor beats one probability for every test by more than 0.1 in mean log score in either split direction, and no permutation p-value is below 0.05. & Plan's expectation: likely. & held \\
10 & P8 & The survey finds at least one dataset with per-test outcomes and times for cells in direct contact, and none whose cell this study's parameters cover without an invented value. & Plan's expectation: uncertain. & failed \\
11 & P1 & For each of the four tests, $f^*$(C2) under every kinetic draw lies within a factor of two of its posterior-mean value, or within 0.025 of it where that is 0. & Plan's expectation: likely. & failed \\
11 & P2 & For each of the four tests, $f^*$(C1) under every kinetic draw lies within a factor of two of its posterior-mean value, or within 0.025 of it where that is 0. & Plan's expectation: uncertain; the hot spot's onset is local and the kinetics set it. & held \\
11 & P3 & Under each of the eight draws, the range of the C2 threshold energy across the four tests exceeds the largest within-test band of that energy across the draws. & Plan's expectation: likely. & held \\
11 & P4 & In the four tests with zero peak-matching fraction, the far face with no delivery exceeds the measured maximum under at least three of four thermal draws and at both contacts. & Plan's expectation: uncertain; the low contact and the smallest radial conductivity are the two that could remove it; the predictions file's analysis: “P4, which asks that it do so at both contacts, fails before any step 11 run” & failed, foreseen \\
11 & P5 & In each of the four tests the band of $f^*$(C2) across the four thermal draws is narrower than its band across the four kinetic draws. & Plan's expectation: uncertain. & failed \\
12 & P1 & With $E_\mathrm{b}$ capped at 33.2 kJ, conduction and radiation alone meet C1 in none of the eight tests at the central contact. & Plan's expectation: likely (step 5 found no ignition at 33.2 kJ with the two-cell study's trigger); the predictions file records it as fixed by step 6's records (Section~\ref{sec:s-registry}). & held \\
12 & P2 & At the high contact the same holds. & Plan's expectation: uncertain; the predictions file records it as close to fixed (step 6, with other tabs; Section~\ref{sec:s-registry}). & held \\
12 & P3 & With $E_\mathrm{e}$ fixed at the Databank's median, the delivered energies $E^*$ across the nine propagated tests still span at least tenfold. & Plan's expectation: likely; the predictions file records it as fixed by step 9's records. & held \\
12 & P4 & With $E_\mathrm{e}$ fixed, the propagated and non-propagated tests' fractions still overlap (the largest $f_\mathrm{max}$ above the smallest $f_\mathrm{t}$(C2)). & Plan's expectation: likely; the predictions file records it as fixed by step 9's records. & held \\
13 & P1 & For PT-Exp1 with a 5 s triangular delivery over the facing half, the resolved $f^*$(C1) lies within 10 \% of the skin model's closed form (0.07701). & likely & held \\
13 & P2 & At 40 s (PT-Exp1, facing half), the cross-section conduction model (0.21775) lies within 5 \% of the resolved bracket's midpoint. & uncertain & held \\
13 & P3 & At 120 s (PT-Exp1, facing half), the resolved threshold lies more than 5 \% below the cross-section conduction model's (0.41450): local ignition. & uncertain & failed \\
13 & P4 & For ST-Exp3 and PT-Exp3, the cross-section model with the trigger's conduction lies within 10 \% of the resolved midpoints at 10 and 60 s, and gives the direction of change from 21.39 s in all four cases. & uncertain & held \\
13 & P7 & The ignition temperature of the posterior-mean kinetics lies between 215 and 240\degC{} (known when registered, so not scored), and the central 95 \% of the 200 draws spans less than 40 K (scored). & likely for the mean, uncertain for the span & held \\
13 & P8 & The two-body model's whole-cell threshold lies within the larger of 10 \% and 300 J of the lumped model's midpoint in at least 28 of the 32 low and high contact searches. & uncertain & held \\
13 & P9 & For PT-Exp1, the ratio $f^*$(C2)/$f^*$(C1) falls from 21.39 to 120 to 360 s and is below 1.5 at 360 s. & likely for the fall, uncertain for the level & held \\
13 & P10 & For PP-Exp3, ST-Exp2 and PT-Exp3, the lumped model's needed delivery falls by 150 to 300 J per kJ of the trigger's body energy across the 4 kJ between shifted runs. & likely for the first two, uncertain for PT-Exp3 & failed \\
14 & P1 & The upper bound of radiation from a cap-end vent to the facing half is below 2 \% of 76.74 kJ. & the plan's line for it says likely; the predictions file's analysis, written before any bound, says P1 is expected to fail & failed, foreseen \\
14 & P2 & The upper bounds of the cap-end paths that are determinable sum to less than 5.96 kJ, the largest need among the seven matched times. & uncertain; the predictions file's analysis says P2 is expected to fail & failed, foreseen \\
14 & P3 & They sum to more than 0.26 kJ, the smallest need. & likely & held \\
14 & P4 & The upper bound of a breach's jet facing the neighbour exceeds 5.96 kJ delivered within 21.39 s. & likely & held \\
14 & P5 & At least one of the six paths is not determinable from published inputs for this cell. & likely & held \\
14 & P6 & For the parallel tests the circuit's upper bound on the neighbour's own heating exceeds 1 kJ. & uncertain; the predictions file's analysis says P6 is expected to hold & held \\
15 & P1 & The reactive cross-section's $E^*$(C8) lies within 15 \% of the midpoint of the three-dimensional bracket in at least three of four tests. & uncertain; the reduced model's values are the searches' centres & held \\
15 & P2 & For PT-Exp4 with the 21.39 s delivery, $E^*$(C8) lies between 6 and 10 kJ. & likely; the reduced model's value, registered with it: 7,062 J & held \\
15 & P3 & In at least three of the four tests $E^*$(C8) is at least 1.25 times the whole-cell clock's threshold $E^*$(C2). & likely for the three strong triggers, uncertain for ST-Exp3; the reduced model's values: 1.49, 1.72, above zero in SP-Exp3, 0.85 & held \\
15 & P4 & For PT-Exp4, in the run at the upper end of the bracket, the mean crosses 255.7\degC{} between 150 and 400 s after the trigger's start. & likely; the reduced model's value: 310.6 s at the upper end of its own bracket & held \\
15 & P5 & In that run the far face reads more than 15 K above 197.3\degC, the reading from which the measured neighbour jumped, when its own jump begins. & likely; the reduced model's value: 240.1\,\textdegree C & held \\
15 & P6 & The constant power after the 5.0 kJ pulse that brings the cell's own runaway to 271.5 s lies between 5 and 25 W. & likely; the reduced model's value: 12.5 W & held \\
15 & P7 & Under none of four kinetic draws and not with the aged SEI does the resolved cell run away at PT-Exp4's 5,045 J, and the least shift of the rate laws that makes it run away is 10 K or more. & likely; the reduced model's values: no runaway in any of the five runs, the least shift 20.0 K & held \\
15 & P8 & PT-Exp4's run with 6,209 J over 240 s, carried to 900 s, runs away. & likely; the reduced model's value, registered with it: no runaway, the mean peaking at 220.5\,\textdegree C & failed \\
\bottomrule
\end{longtable}
}
}{}

\section{The thresholds and needs of every test}
\label{sec:s-thresholds}

Table~\ref{tab:s-thresholds} gives, test by test, the brackets that Fig.~2 and Fig.~4a of the main text draw. The local and lumped brackets of most tests come from searches over fractions of each test's ejected energy that halved the bracket to 0.0125 of it; four tests' brackets were narrowed further with one ejected energy for every test, and the lumped brackets of PT-Exp4 and ST-Exp3 by the search for the measured time, which ended at the threshold. The cross-section's brackets are at most 1\,\% of their midpoints; the three-dimensional searches started 5\,\% either side of the cross-section's value.

\IfFileExists{tables/s3_thresholds.tex}{{\footnotesize\setlength{\tabcolsep}{3pt}
\begin{longtable}{p{0.1\textwidth}p{0.13\textwidth}p{0.15\textwidth}p{0.15\textwidth}p{0.15\textwidth}p{0.18\textwidth}}
\caption{The thresholds of the sixteen tests with the 21.4~s pulse (kJ; a bracket is the largest energy run that does not meet the criterion and the smallest that does), and, for the nine propagated tests, their needs: the least energy with which the lumped neighbour and the resolved neighbour (by the reactive cross-section) run away within 10~s of the measured time, counted as delivered by the run's own runaway. An energy marked $^{t}$ is the least energy run whose runaway comes before the measured time when the next lower run, 0.5\,\% of the test's ejected energy below, gives none within 900~s: the measured time is reproduced only between the two, within that resolution of the threshold, and the value, the higher, can exceed the bracket beside it. Supplementary Table~\ref{tab:s-needs} gives the other delivery shapes.}\label{tab:s-thresholds}\\
\toprule
Test & Hottest point & Lumped & Resolved, cross-section & Resolved, three-dimensional & Needs: lumped; resolved \\
\midrule
\endfirsthead
\multicolumn{6}{l}{\textit{Continued}} \\
\toprule
Test & Hottest point & Lumped & Resolved, cross-section & Resolved, three-dimensional & Needs: lumped; resolved \\
\midrule
\endhead
PT-Exp1 & (4.36, 4.53] & (12.85, 12.89] & (10.75, 10.81] &  &  \\
PT-Exp2 & (0.00, 1.68] & (1.68, 3.37] & (6.33, 6.38] &  &  \\
PT-Exp3 & (2.61, 2.68] & (7.36, 8.58] & (8.56, 8.62] &  &  \\
PT-Exp4 & (0.00, 1.55] & (4.66, 4.85] & (7.04, 7.08] & (6.71, 7.42] & 4.85$^{t}$; 7.44$^{t}$ \\
PP-Exp1 & 0 & 0 & (4.53, 4.56] &  & 2.97; 6.04 \\
PP-Exp2 & (0.00, 1.59] & (3.18, 4.77] & (6.83, 6.88] & (6.51, 7.20] & 5.96; 7.65 \\
PP-Exp3 & (0.00, 0.85] & (2.56, 3.41] & (6.21, 6.25] &  &  \\
PP-Exp4 & 0 & (0.00, 1.82] & (5.58, 5.62] &  & 1.82; 6.05$^{t}$ \\
ST-Exp1 & 0 & (0.00, 1.98] & (4.42, 4.44] &  &  \\
ST-Exp2 & 0 & (1.80, 2.70] & (5.58, 5.62] &  &  \\
ST-Exp3 & (4.04, 4.20] & (11.95, 12.11] & (10.25, 10.31] & (9.77, 10.80] & 12.11$^{t}$; 10.31$^{t}$ \\
ST-Exp4 & 0 & (0.00, 1.89] & (5.03, 5.06] &  & 1.18; 5.51$^{t}$ \\
SP-Exp1 & 0 & (1.82, 3.63] & (5.50, 5.54] &  & 2.73; 5.96 \\
SP-Exp2 & 0 & (0.00, 1.89] & (5.06, 5.08] &  & 3.78; 6.50 \\
SP-Exp3 & 0 & 0 & (3.52, 3.54] & (3.70, 4.06] & 0.26; 4.58 \\
SP-Exp4 & (0.00, 1.54] & (4.61, 6.15] & (7.08, 7.12] &  &  \\
\bottomrule
\end{longtable}
}
}{}

\section{The checks}
\label{sec:s-checks}

Every run's energy account closes: the heat taken in through the boundary paths, the reaction heat and the delivered energy equal the change in stored heat to the residuals of Table~\ref{tab:s-checks}, against an acceptance of $10^{-9}$ of the heat taken in, and every delivery equals its pulse's integral up to the run's end. Each step that changed the configuration of the resolved neighbour halved every mesh spacing, and separately every step limit, on a run near its threshold. Every registered number and verdict was recomputed by separate code and read again, before the step's texts were written, by a check that had not drafted them.

\begin{table}[!htbp]
\centering
\footnotesize
\caption{The checks of the steps whose runs the main text reports: the largest energy residual of each step's runs, as a fraction of the heat taken in, and the change in the hottest point's peak rise (or a runaway time) with the mesh spacing halved and with every step limit halved.}
\label{tab:s-checks}
\setlength{\tabcolsep}{4pt}
\begin{tabular}{rrlll}
\toprule
Step & Runs & Largest energy residual & Mesh halved & Step limits halved \\
\midrule
5 & 5,248 & $8.5\times10^{-12}$ & 0.27 and 0.31\,\% & 0.0035 and 0.0040\,\% \\
6 & 506 & $8.5\times10^{-12}$ & 0.28 and 0.32\,\% & 0.0033 and 0.0037\,\% \\
7 & 186 & $4.34\times10^{-12}$ & 0.15\,\% & 0.0008\,\% \\
8 & 230 & $1.97\times10^{-11}$ & 0.52\,\% & 0.0021\,\% \\
9 & 670 & $6.8\times10^{-12}$ & 0.32\,\% & 0.0006\,\% \\
10 & 30 & $3.3\times10^{-12}$ & \multicolumn{2}{l}{the configuration of step 9, unchanged} \\
11 & 216 & $6.1\times10^{-12}$ & \multicolumn{2}{l}{the configuration of step 9, unchanged} \\
12 & 10 & $1.2\times10^{-12}$ & \multicolumn{2}{l}{the configuration of step 9, unchanged} \\
13 & 389 & $2.99\times10^{-12}$ & \multicolumn{2}{l}{the configuration of step 9, unchanged} \\
15 & 27 & $1.97\times10^{-12}$ & 1.5\,\% (runaway time) & 0.01\,\% (runaway time) \\
\bottomrule
\end{tabular}
\end{table}

\section{The inputs of the bounds on what each path can deliver}
\label{sec:s-inputs}

Each input of the bounds of Table~3 of the main text was read at its source, with the page used, and labelled as measured on this cell, in this rig, on another cell, or as a relation; a value measured on another cell widens a range and sets no value for this one (Table~\ref{tab:s-inputs}). Where no source gives a fraction, it enters at 1, the largest it can be, so each upper bound is what the published inputs cannot be made to exceed. Radiation was bounded by tracing four to eight million rays per point from emitters placed on a grid beyond the trigger's end: on the vent's axis the largest share of an emitter's rays that strike the facing half is about 0.0037, and in front of the trigger's end face, near the line of contact, the share rises toward a quarter. With the holder able to reflect, the geometry stops mattering, and radiation is bounded by a black body at the flame temperature filling the facing half's view.

\begin{table}[!htbp]
\centering
\footnotesize
\caption{The inputs of the bounds. Whose: measured on this cell (the LG M50), in this rig, on another cell, or a relation.}
\label{tab:s-inputs}
\begin{tabular}{p{0.31\textwidth}p{0.25\textwidth}p{0.2\textwidth}p{0.11\textwidth}}
\toprule
Input & Value & Source & Whose \\
\midrule
Largest energy in the ejecta, ten nail tests & 89.34~kJ (median 76.74~kJ) & \cite{finegan2024databank} & this cell \\
Largest heat of burning in open air & 136.5~kJ & \cite{vendra2022numerical}, Table 2 & another cell \\
Largest flame temperature & 1500\degC, a battery module & cited in \cite{chen2020experimental}, p.~2 & another cell \\
Largest vent gas velocity & 398.5~m\,s$^{-1}$ & \cite{gillich2024mechanical} & another cell \\
Shortest runaway or venting & 1.131~s & \cite{gillich2024mechanical} & another cell \\
Longest combustion of a 21700 cell & about 300~s & \cite{chen2023comprehensive}, p.~11 & another cell \\
Sidewall ruptures of the trigger & 3 of the 16 tests & \cite{fransson2024exploring}, Tables 3 and 4 & this cell \\
Sidewall breaches in the Databank's nail tests & none of 10 & \cite{finegan2024databank} & this cell \\
A rupture's jet fire on neighbours & neighbours ran away 2 to 5~s later & \cite{read2024sidewall}, p.~5 & this cell \\
Holder & aluminium frame with 3D-printed inserts, about 2~mm of contact; material not given & \cite{pfaff2023insitu}, p.~194 & this rig \\
Discharge of the neighbour through the trigger & stated for the parallel tests & \cite{fransson2024exploring} & this rig \\
Stored energy at 100\,\% state of charge & 65.52~kJ (18.2~Wh) & \cite{finegan2024databank} & this cell \\
Resistance's rise as the cell cools & about five times from 45 to 10\degC & \cite{popp2020antemortem}, p.~172 & this cell \\
Convection from a cylinder in cross flow & Churchill and Bernstein's relation & \cite{churchill1977correlating,lienhard2024heat} & relation \\
Properties of hydrogen and air at 1 atm & tabulated to 1000 and 1500~K & \cite[Table A.6]{lienhard2024heat} & relation \\
Radiant fraction, absorptivity, share of gas passing, share carried by a breach & no source; each enters at 1 & & limit \\
\bottomrule
\end{tabular}
\end{table}

The neighbour's direct-current resistance, $30\pm6$~m$\Omega$ at 25\degC, is the manufacturer's specification of the cell (LG Chem, Product Specification, Lithium Ion INR21700 M50 18.2~Wh, read in November 2026); it sets only the rate at which the discharge path could deliver, 15.6~kJ within 21.4~s if the resistance kept its value at 25\degC, while the bound in Table~3 of the main text is the whole stored energy.

\section{The one published neighbour history}
\label{sec:s-history}

No raw temperature series of the two-cell tests is published. The supplement of Fransson et al.~\cite{fransson2024exploring} draws the thermocouple records of two tests in its Figure~10, PT-Exp4, which propagated, and PP-Exp3, which did not. The figure is a vector drawing, so the series were read out of the curves' paths rather than off pixels, twice, by separate code and separate calibrations, and resampled every 0.5~s from 0 to 500~s. The two readings differ by at most 0.0005~K on PT-Exp4's neighbour curve; the figure's own resolution, half the curves' stroke width, is 2.47~K and 0.45~s, and the recorded uncertainty is the larger. A third reading, from the page rendered at 1200~dpi, lies within the recorded uncertainty plus one pixel of the vector reading. PT-Exp4's record: 32.2\degC{} at 0~s; a spike to 102.0\degC{} at 3.1~s as the trigger runs away, back to 79.9\degC{} at 8.4~s; 141.0\degC{} at 45~s; a slow rise to 197.3\degC{} at 271.50~s, where the figure marks the onset (Table~3 of the paper gives 272~s); then the jump. The main text takes the opening 8.5~s as the thermocouple's exposure at the trigger's runaway rather than the can's temperature.

\section{Other datasets}
\label{sec:s-survey}

A survey of published datasets with per-test outcomes and times for cell-to-cell propagation asked whether the model's inputs cover their cell without an invented value (Table~\ref{tab:s-survey}). One other dataset has a cell they cover: Read, Jones and Marco~\cite{read2024sidewall}, seven LG M50 cells 3~mm apart, whose propagated tests followed the trigger's sidewall rupture, a path this model does not represent.

\begin{table}[!htbp]
\centering
\footnotesize
\caption{The datasets the survey found.}
\label{tab:s-survey}
\begin{tabular}{p{0.3\textwidth}p{0.21\textwidth}p{0.17\textwidth}p{0.17\textwidth}}
\toprule
Dataset & Cell & Arrangement & The model's inputs cover the cell \\
\midrule
Fransson et al. 2024~\cite{fransson2024exploring}, M50 tests & LG M50 & two cells in contact & yes \\
Fransson et al. 2024, P42A tests & Molicel P42A & two cells in contact & no \\
Read, Jones and Marco 2024~\cite{read2024sidewall} & LG M50 & seven cells, 3~mm gaps & yes \\
Read et al. 2024~\cite{read2024barrier} & unnamed, M50 specification & 34 cells, 1~mm gaps & not established \\
Allen et al. 2025~\cite{allen2025sleeves} & Molicel P45B & seven cells, 3~mm & no \\
Lamb et al. 2015~\cite{lamb2015failure} & LiCoO$_2$, 18650 and pouch & cells in contact & no \\
Said et al. 2019~\cite{said2019cascading} & LiCoO$_2$, 18650 & arrays & no \\
Hoelle et al. 2023~\cite{hoelle2023simulation} & prismatic & five-cell stack & no \\
Sch\"oberl et al. 2024~\cite{schoberl2024propagation} & prismatic NMC and LFP & braced stack & no \\
Torres-Castro et al. 2020~\cite{torrescastro2020passive} & pouch & stacks & no \\
\bottomrule
\end{tabular}
\end{table}

\bibliographystyle{elsarticle-num}
\bibliography{references}

\end{document}
